\section{The Euclidean Geometry Available Before Coordinates}

Before introducing coordinates, we recall the geometric world that is already
available. At this stage there are no ordered pairs, no axes, and no vectors.
There are points, straight lines, segments, circles, and angles, together with
relations such as incidence, betweenness, congruence, perpendicularity, and
parallelism.

\begin{remarkbox}
Coordinates will later translate this geometry into arithmetic. They do not
create the geometry. We can already construct and compare geometric objects
before assigning numbers to all points and angles.
\end{remarkbox}

\subsection{Three Different Levels Must Not Be Confused}

We shall distinguish three kinds of statements.

\begin{enumerate}
    \item \emph{Euclid's postulates} describe the geometric world being assumed.
    \item \emph{Straightedge-and-compass rules} describe which elementary
          drawings and intersections may be used in a construction.
    \item \emph{Derived constructions} must be built from those elementary
          operations and justified by a proof.
\end{enumerate}

In particular, a perpendicular bisector, a perpendicular through a point, an
angle bisector, or a parallel through a point will not be treated as primitive
construction postulates. They will be constructed.

\subsection{Euclid's Five Postulates}

The following are the five postulates used in the first book of Euclid's
\emph{Elements}. The diagrams illustrate their meaning; they are not proofs.

\subsubsection{Postulate I: Joining Two Points}

A straight segment can be drawn from any point to any other point.

\begin{center}
\begin{tikzpicture}[scale=0.9]
    \coordinate (A) at (-2.4,-0.6);
    \coordinate (B) at (2.1,0.8);
    \draw[subset] (A) -- (B);
    \fill[geometricSubsetColor] (A) circle (2pt) node[below] {$A$};
    \fill[geometricSubsetColor] (B) circle (2pt) node[above] {$B$};
    \node at (0,-1.55) {a segment is drawn from $A$ to $B$};
\end{tikzpicture}
\end{center}

\subsubsection{Postulate II: Extending a Segment}

A finite straight segment can be extended continuously in the same straight
line.

\begin{center}
\begin{tikzpicture}[scale=0.9]
    \coordinate (A) at (-2.2,0);
    \coordinate (B) at (1.2,0);
    \coordinate (C) at (3.1,0);
    \draw[subset] (A) -- (B);
    \draw[helper,->] (B) -- (C);
    \fill[geometricSubsetColor] (A) circle (2pt) node[below] {$A$};
    \fill[geometricSubsetColor] (B) circle (2pt) node[below] {$B$};
    \node[below] at (C) {$C$};
    \node at (0,-1.15) {$AB$ is extended beyond $B$};
\end{tikzpicture}
\end{center}

\subsubsection{Postulate III: Drawing a Circle}

A circle can be drawn with any chosen centre and any chosen radius.

\begin{center}
\begin{tikzpicture}[scale=0.9]
    \coordinate (O) at (0,0);
    \coordinate (A) at (2.2,0);
    \draw[helper] (O) circle (2.2);
    \draw[subset] (O) -- (A) node[midway,below] {$r$};
    \fill[geometricSubsetColor] (O) circle (2pt) node[below left] {$O$};
    \fill[geometricSubsetColor] (A) circle (2pt) node[right] {$A$};
    \node at (0,-2.75) {centre $O$ and radius $r$ determine a circle};
\end{tikzpicture}
\end{center}

This postulate is what permits a length to be copied geometrically without
first assigning a number to it.

\subsubsection{Postulate IV: Equality of Right Angles}

All right angles are congruent.

\begin{center}
\begin{tikzpicture}[scale=0.85]
    \draw[subset] (-3.2,0) -- (-0.2,0);
    \draw[subset] (-1.7,-1.3) -- (-1.7,1.6);
    \draw (-1.42,0) -- (-1.42,0.28) -- (-1.7,0.28);

    \draw[subset] (0.8,-0.8) -- (3.5,0.7);
    \draw[subset] (2.9,-1.25) -- (1.4,1.45);
    \draw (2.24,0.02) -- (2.10,0.27) -- (1.85,0.13);

    \node at (0,-2.0) {every right angle represents the same geometric magnitude};
\end{tikzpicture}
\end{center}

This postulate does not say that a perpendicular may simply be drawn without
construction. It says that once right angles have been constructed, they are
all equal.

\subsubsection{Postulate V: The Parallel Postulate}

If a transversal meets two straight lines and the interior angles on one side
sum to less than two right angles, then the two lines, when extended, meet on
that side.

\begin{center}
\begin{tikzpicture}[scale=0.85]
    \draw[subset] (-3.4,1.15) -- (3.5,0.15) node[right] {$\ell$};
    \draw[subset] (-3.4,-1.2) -- (3.5,-0.35) node[right] {$m$};
    \draw[helper] (-1.0,-2.0) -- (0.6,2.0) node[above] {$t$};
    \node at (0,-2.45) {the angle condition determines on which side $\ell$ and $m$ meet};
\end{tikzpicture}
\end{center}

A familiar modern equivalent is Playfair's formulation: through a point outside
a given line there passes exactly one parallel to that line. We shall use that
form later, after the relation between the two formulations has been made clear.

\subsection{Elementary Rules for Straightedge-and-Compass Work}

For explicit constructions we use the following operational rules.

\begin{enumerate}
    \item Through two already constructed distinct points, draw the straight line.
    \item With an already constructed point as centre and an already constructed
          segment as radius, draw a circle.
    \item Retain intersection points of two constructed lines, a line and a
          circle, or two circles whenever such intersections exist.
\end{enumerate}

\begin{center}
\begin{tikzpicture}[scale=0.9]
    \coordinate (A) at (-2.3,-0.4);
    \coordinate (B) at (2.1,0.9);
    \draw[subset] (A) -- (B);
    \fill[geometricSubsetColor] (A) circle (2pt) node[below] {$A$};
    \fill[geometricSubsetColor] (B) circle (2pt) node[above] {$B$};
    \node at (0,-1.35) {draw the line through two constructed points};
\end{tikzpicture}
\end{center}

\begin{center}
\begin{tikzpicture}[scale=0.9]
    \coordinate (O) at (0,0);
    \coordinate (R) at (2.0,0);
    \draw[helper] (O) circle (2.0);
    \draw[subset] (O) -- (R) node[midway,below] {$r$};
    \fill[geometricSubsetColor] (O) circle (2pt) node[below left] {$O$};
    \fill[geometricSubsetColor] (R) circle (2pt);
    \node at (0,-2.55) {draw a circle from a constructed centre and radius};
\end{tikzpicture}
\end{center}

\begin{center}
\begin{tikzpicture}[scale=0.85]
    \coordinate (C) at (-1.5,0);
    \coordinate (D) at (1.5,0);
    \draw[helper] (C) circle (2.15);
    \draw[helper] (D) circle (2.15);
    \fill (C) circle (1.7pt) node[left] {$C$};
    \fill (D) circle (1.7pt) node[right] {$D$};
    \fill[geometricSubsetColor] (0,1.55) circle (2pt) node[above] {$P$};
    \fill[geometricSubsetColor] (0,-1.55) circle (2pt) node[below] {$Q$};
    \node at (0,-2.65) {retain the intersection points as new constructible points};
\end{tikzpicture}
\end{center}

These rules are the actual construction toolkit. The constructions below must
be reduced to repeated uses of them.

\subsection{Constructions Derived from the Toolkit}

\subsubsection{The Perpendicular Bisector and the Midpoint}

Let $A$ and $B$ be distinct points. Draw two circles with the same radius,
one centred at $A$ and one centred at $B$, choosing the radius larger than half
of $AB$. Let their intersection points be $P$ and $Q$.

Because $P$ lies on both circles,
\[
PA=PB,
\]
and similarly,
\[
QA=QB.
\]
Thus both $P$ and $Q$ are equidistant from $A$ and $B$. The line $PQ$ is the
perpendicular bisector of $AB$.

\begin{center}
\begin{tikzpicture}[scale=0.85]
    \coordinate (A) at (-2.2,0);
    \coordinate (B) at (2.2,0);
    \coordinate (P) at (0,2.4);
    \coordinate (Q) at (0,-2.4);
    \draw[helper] (A) circle (3.25);
    \draw[helper] (B) circle (3.25);
    \draw[subset] (A) -- (B);
    \draw[very thick] (0,-3.0) -- (0,3.0);
    \fill[geometricSubsetColor] (A) circle (2pt) node[below] {$A$};
    \fill[geometricSubsetColor] (B) circle (2pt) node[below] {$B$};
    \fill[geometricSubsetColor] (P) circle (2pt) node[above] {$P$};
    \fill[geometricSubsetColor] (Q) circle (2pt) node[below] {$Q$};
\end{tikzpicture}
\end{center}

The intersection of $PQ$ with $AB$ is the midpoint of $AB$. Thus midpoint and
perpendicularity have been constructed from equal radii and intersections.

\subsubsection{A Perpendicular Through a Point on a Line}

Let $P$ lie on a line $\ell$. Choose points $A$ and $B$ on $\ell$ with
\[
PA=PB.
\]
Construct the perpendicular bisector of $AB$. Since $P$ is the midpoint of
$AB$, this perpendicular bisector passes through $P$ and is perpendicular to
$\ell$.

\begin{center}
\begin{tikzpicture}[scale=0.9]
    \coordinate (A) at (-2.4,0);
    \coordinate (P) at (0,0);
    \coordinate (B) at (2.4,0);
    \draw[subset] (-3.4,0) -- (3.4,0) node[right] {$\ell$};
    \draw[very thick] (0,-2.0) -- (0,2.4) node[above] {$m$};
    \draw (0.28,0) -- (0.28,0.28) -- (0,0.28);
    \fill[geometricSubsetColor] (A) circle (2pt) node[below] {$A$};
    \fill[geometricSubsetColor] (P) circle (2pt) node[below right] {$P$};
    \fill[geometricSubsetColor] (B) circle (2pt) node[below] {$B$};
    \node at (0,-2.45) {$m$ is constructed through $P$ and satisfies $m\perp\ell$};
\end{tikzpicture}
\end{center}

\subsubsection{A Perpendicular Through a Point Outside a Line}

Let $P$ lie outside a line $\ell$. Draw a circle centred at $P$ large enough
to meet $\ell$ at two points $A$ and $B$. Then
\[
PA=PB.
\]
Construct the perpendicular bisector of $AB$. Because $P$ is equidistant from
$A$ and $B$, it lies on that perpendicular bisector. Hence the constructed
line passes through $P$ and is perpendicular to $\ell$.

\begin{center}
\begin{tikzpicture}[scale=0.9]
    \coordinate (P) at (0,2.0);
    \coordinate (A) at (-2.1,0);
    \coordinate (B) at (2.1,0);
    \draw[subset] (-3.5,0) -- (3.5,0) node[right] {$\ell$};
    \draw[helper] (P) circle (2.9);
    \draw[very thick] (0,-1.0) -- (0,2.8) node[above] {$m$};
    \draw (0.28,0) -- (0.28,0.28) -- (0,0.28);
    \fill[geometricSubsetColor] (P) circle (2pt) node[right] {$P$};
    \fill[geometricSubsetColor] (A) circle (2pt) node[below] {$A$};
    \fill[geometricSubsetColor] (B) circle (2pt) node[below] {$B$};
\end{tikzpicture}
\end{center}

The existence of a perpendicular has therefore been proved by construction; it
was not inserted as an additional postulate.

\subsubsection{Bisecting an Angle}

Let two rays with common endpoint $O$ form an angle. Draw a circle centred at
$O$ meeting the rays at $A$ and $B$. Then draw equal-radius circles centred at
$A$ and $B$, and let one of their intersections inside the angle be $P$.
Since
\[
OA=OB,
\qquad
AP=BP,
\qquad
OP=OP,
\]
the triangles $OAP$ and $OBP$ are congruent. Therefore
\[
\angle AOP=\angle POB.
\]
The ray $OP$ is the angle bisector.

\subsubsection{Constructing a Parallel Through a Point}

Let $P$ lie outside a line $\ell$. First construct a line $m$ through $P$
perpendicular to $\ell$. Then construct a line $n$ through $P$ perpendicular
to $m$. Since $n$ and $\ell$ are perpendicular to the same line, they have the
same direction and do not meet:
\[
n\parallel\ell.
\]
The fifth postulate guarantees the uniqueness of this parallel in Euclidean
geometry.

\subsection{Parallel Lines and Proportional Segments}

A central fact of Euclidean geometry is that parallel lines preserve ratios.
This is the content of Thales' theorem in the form needed here.

\begin{theorembox}
Let $D$ lie on $AB$ and $E$ lie on $AC$ in a triangle $ABC$. If
\[
DE\parallel BC,
\]
then
\[
\frac{AD}{AB}
=
\frac{AE}{AC}
=
\frac{DE}{BC}.
\]
\end{theorembox}

\begin{center}
\begin{tikzpicture}[scale=0.9]
    \coordinate (A) at (0,3.6);
    \coordinate (B) at (-3.2,0);
    \coordinate (C) at (3.8,0);
    \coordinate (D) at (-1.6,1.8);
    \coordinate (E) at (1.9,1.8);
    \draw[subset] (A) -- (B) -- (C) -- cycle;
    \draw[very thick] (D) -- (E);
    \fill[geometricSubsetColor] (A) circle (2pt) node[above] {$A$};
    \fill[geometricSubsetColor] (B) circle (2pt) node[below left] {$B$};
    \fill[geometricSubsetColor] (C) circle (2pt) node[below right] {$C$};
    \fill[geometricSubsetColor] (D) circle (2pt) node[left] {$D$};
    \fill[geometricSubsetColor] (E) circle (2pt) node[right] {$E$};
    \node at (0.15,2.05) {$DE\parallel BC$};
\end{tikzpicture}
\end{center}

The smaller triangle $ADE$ and the larger triangle $ABC$ have the same three
angles. Their corresponding sides differ by one common scale factor. This is
the basic phenomenon of similarity.

\begin{definitionbox}
Two triangles are similar when their corresponding angles are equal. In that
case their corresponding side lengths are proportional.
\end{definitionbox}

Similarity lets us change the size of a figure while preserving its shape. It
is the reason ratios of corresponding lengths can depend only on angles.

\subsection{Sine and Cosine Before Coordinates}

Take a right triangle containing an acute angle $\alpha$. Let $h$ denote the
hypotenuse, let $a$ be the leg adjacent to $\alpha$, and let $b$ be the leg
opposite $\alpha$.

\begin{center}
\begin{tikzpicture}[scale=0.95]
    \coordinate (A) at (0,0);
    \coordinate (B) at (4.8,0);
    \coordinate (C) at (4.8,2.9);
    \draw[subset] (A) -- (B) -- (C) -- cycle;
    \draw (B) ++(-0.3,0) -- ++(0,0.3) -- ++(0.3,0);
    \draw (0.85,0) arc[start angle=0,end angle=31.1,radius=0.85];
    \node at (1.08,0.28) {$\alpha$};
    \node[below] at (2.4,0) {$a$};
    \node[right] at (4.8,1.45) {$b$};
    \node[above left] at (2.4,1.45) {$h$};
\end{tikzpicture}
\end{center}

We define
\[
\cos\alpha=\frac{a}{h},
\qquad
\sin\alpha=\frac{b}{h}.
\]

These expressions are well defined. Indeed, any two right triangles containing
the same acute angle $\alpha$ have the same third angle and are therefore
similar. Their corresponding sides are proportional, so the ratios $a/h$ and
$b/h$ do not depend on the particular triangle chosen. They depend only on the
angle.

This is the essential role of Thales' theorem and similarity: they turn the two
ratios above into genuine functions of an angle.

The Pythagorean theorem gives
\[
a^2+b^2=h^2.
\]
Dividing by $h^2$ yields
\[
\left(\frac{a}{h}\right)^2+
\left(\frac{b}{h}\right)^2=1,
\]
and hence
\[
\cos^2\alpha+\sin^2\alpha=1.
\]
Thus the fundamental trigonometric identity is a direct reformulation of the
Pythagorean theorem.

\subsection{Angles and Their Numerical Measure}

An angle exists geometrically before it is assigned a number. We can compare,
copy, add, and bisect angles without degrees or radians. A numerical angle
measure requires an additional choice of unit.

Degrees divide one full turn into $360$ equal parts. Radians use the geometry
of a circle. If an angle at the centre of a circle cuts off an arc of length
$s$ on a circle of radius $r$, its radian measure is
\[
\theta=\frac{s}{r}.
\]
The ratio does not depend on the size of the circle: scaling the circle by a
factor scales both $s$ and $r$ by that factor. Again, similarity makes the
quantity well defined.

\begin{remarkbox}
At this point no vector has been introduced. We have only Euclidean objects,
postulates, constructions, similarity, trigonometric ratios, and numerical
angle measure. Vectors will appear only in the next chapter, after the
Cartesian coordinate system has been constructed.
\end{remarkbox}

The next step is to choose two perpendicular lines, an origin, directions, and
a unit length. Those choices will convert the Euclidean plane into a Cartesian
coordinate plane.
